%% capitulos/fundamentacao.tex — Capítulo 2: Fundamentação Teórica
%%
%% Bloco 1 (Conteúdo Teórico) escrito por Claude com base em buscas
%% bibliográficas no scite (DOIs verificados). Sinalizado em laranja itálico
%% pelo ambiente claudedraft. Trave a seção quando o conteúdo for aprovado:
%%   python -m src.writer --lock fundamentacao/conteudo

\chapter{Fundamentação Teórica}
\label{ch:fundamentacao}

% WRITER:BEGIN:fundamentacao/conteudo
\begin{claudedraft}

Este capítulo apresenta o aparato conceitual e operacional sobre o qual
se constrói a metodologia da dissertação. A
\autoref{sec:teoria-informacao} retoma a entropia de Shannon e seus
desdobramentos em entropia diferencial e divergências, dando base
formal às medidas que vêm a seguir. A \autoref{sec:apen-sampen}
percorre a transição da Entropia Aproximada (ApEn) para a Entropia de
Amostra (SampEn), explicitando os vieses que motivaram o refinamento e
a interpretação probabilística de SampEn como log-verossimilhança
condicional. As \autoref{sec:sampen2d} e \autoref{sec:mse} estendem
SampEn ao domínio espacial bidimensional e ao plano multiescala,
respectivamente, fornecendo as duas leituras complementares sobre a
matriz de pesos de uma camada convolucional. Por fim, a
\autoref{sec:lmc} apresenta a complexidade estatística de López-Ruiz,
Mancini e Calbet (LMC), discute a controvérsia sobre normalização e
intensividade que se seguiu à proposta original e justifica a
convenção $C = H \times D$ adotada no restante do trabalho.

\section{Sistemas complexos, auto-organização e emergência}
\label{sec:sistemas-complexos}

Esta seção fixa o vocabulário conceitual sobre o qual as medidas das seções
seguintes se apoiam. A justificativa para sua presença é direta: o restante
desta dissertação afirma que os pesos de uma rede em treinamento apresentam
regimes de \emph{auto-organização}, e essa afirmação não tem conteúdo
verificável enquanto o termo não for definido.

\subsection{A noção de sistema complexo}

Um sistema complexo pode ser informalmente caracterizado como uma rede de
muitos componentes que interagem entre si e cujo comportamento coletivo surge e
evolui por auto-organização \cite{sayama2015}. A caracterização carrega duas
exigências que não são triviais. A primeira é a \emph{multiplicidade} de
componentes interagentes. A segunda, mais restritiva, é que a informação
relevante do sistema resida nas \emph{relações} entre as partes, e não nas
partes isoladamente; se o comportamento do todo for recuperável pela simples
soma dos comportamentos individuais, o sistema é agregado, não complexo.

Uma rede neural profunda em treinamento satisfaz ambas as exigências. Os
milhões de pesos constituem componentes que interagem indiretamente, pela
retropropagação do gradiente, e nenhum peso isolado codifica a função aprendida
--- o conteúdo do modelo está na configuração conjunta. Essa observação não é
retórica: é ela que autoriza aplicar ao treinamento o instrumental
desenvolvido para sistemas complexos, em lugar de tratar a rede apenas como um
problema de otimização de alta dimensão.

\subsection{Emergência}

Denomina-se \emph{emergente} a propriedade que se manifesta no nível
macroscópico do sistema e que não é redutível às propriedades dos componentes
isolados \cite{sayama2015}. O conceito exige, portanto, a distinção explícita de
dois níveis de descrição: o microscópico, dos componentes, e o macroscópico, do
comportamento observável.

Na presente dissertação essa distinção é operacional, e não apenas conceitual.
O nível microscópico é o dos valores individuais dos pesos; o macroscópico é o
da capacidade de generalização do modelo, aferida em dados não vistos. A
hipótese central pode então ser reformulada nesses termos: pretende-se
verificar se uma estatística calculada sobre o nível microscópico --- a
distribuição e a organização dos pesos --- acompanha, e eventualmente antecipa,
a transição de regime observada no nível macroscópico.

\subsection{Auto-organização}

Chama-se \emph{auto-organização} ao processo pelo qual um sistema desenvolve
estrutura macroscópica sem que essa estrutura seja imposta por um controle
externo ou por um plano centralizado \cite{sayama2015}. Emergência e
auto-organização são conceitos distintos, ainda que frequentemente tratados como
sinônimos: a primeira descreve uma relação entre níveis de descrição em um
instante, e a segunda descreve a evolução temporal do sistema.

O treinamento por descida de gradiente é, nessa acepção, um processo de
auto-organização. Nenhum agente externo especifica o valor que cada peso deve
assumir; os valores resultam da aplicação iterada de uma regra local de
atualização, e a estrutura que se observa ao final --- filtros que respondem a
padrões visuais, camadas que se especializam --- não foi projetada peso a peso.

Interessa particularmente a este trabalho a observação de que sistemas em
auto-organização não evoluem de modo homogêneo: eles frequentemente apresentam
\emph{transições de regime}, nas quais o comportamento qualitativo muda em
intervalos relativamente curtos \cite{sayama2015}. É exatamente essa a estrutura
que a hipótese desta dissertação pressupõe. Se o treinamento consistisse em um
processo homogêneo de melhoria gradual, não haveria transição a ser detectada, e
a busca por um marcador do início do sobreajuste careceria de objeto. A
existência de dois regimes sucessivos no aprendizado --- um em que a rede extrai
regularidades generalizáveis e outro em que passa a memorizar exemplos
inconsistentes --- é documentada empiricamente por
\citeonline{arpit2017} e constitui o pressuposto factual da proposta.

\subsection{Delimitação: o que não se pretende afirmar}

Convém delimitar o alcance do enquadramento para evitar leitura excessiva. O
aparato de sistemas complexos é empregado aqui em sentido \emph{descritivo}, e
não como apelo a um mecanismo explicativo. Não se afirma que o treinamento de
uma rede neural seja formalmente equivalente a algum dos modelos canônicos da
área --- autômatos celulares, sistemas de reação-difusão ou redes dinâmicas
\cite{sayama2015} ---, nem se pretende derivar as medidas empregadas a partir de
primeiros princípios da teoria de sistemas complexos.

O que o enquadramento fornece é mais modesto e suficiente: a legitimidade de
tratar a configuração dos pesos como estado de um sistema cuja organização
macroscópica varia no tempo, e cujas transições podem, em princípio, ser
detectadas por estatísticas calculadas sobre esse estado. As medidas
específicas utilizadas provêm da teoria da informação e da análise de séries,
apresentadas nas seções seguintes, e sua validade é estabelecida empiricamente
no Capítulo~\ref{ch:resultados}, não por filiação teórica.

\section{Entropia de Shannon e teoria da informação}
\label{sec:teoria-informacao}

A entropia de Shannon \cite{shannon1948} é a pedra angular sobre a
qual se assenta toda a discussão subsequente. Para uma variável
aleatória discreta $X$ que assume valores no alfabeto finito
$\mathcal{X} = \{x_1, x_2, \ldots, x_N\}$ com distribuição de
probabilidade $P = (p_1, p_2, \ldots, p_N)$, a entropia é definida por
\begin{equation}
  H(X) = -\sum_{i=1}^{N} p_i \log_2 p_i,
  \label{eq:shannon-discrete}
\end{equation}
adotando-se, por convenção, $0 \log_2 0 = 0$. A escolha do logaritmo
na base $2$ implica que $H$ é medida em bits, interpretação
diretamente associada ao número médio mínimo de dígitos binários
necessário para codificar uma realização de $X$. A entropia satisfaz
$0 \le H(X) \le \log_2 N$, com extremos alcançados, respectivamente,
pela distribuição degenerada (existe $i$ tal que $p_i = 1$) e pela
distribuição uniforme ($p_i = 1/N$ para todo $i$). Caracterizações
axiomáticas mostram que essas propriedades, junto da continuidade e da
aditividade, determinam $H$ a menos de um fator multiplicativo
\cite{coverthomas2006}.

Quando $X$ é uma variável contínua com função de densidade de
probabilidade $f$, a generalização natural --- a entropia diferencial
--- é dada por
\begin{equation}
  h(X) = -\int_{\mathcal{X}} f(x) \log_2 f(x)\, dx,
  \label{eq:shannon-diff}
\end{equation}
sempre que a integral exista. Diferentemente da versão discreta,
$h(X)$ pode assumir valores negativos e não é invariante sob mudança
de escala, o que limita seu uso direto como índice de
``incerteza''. Em aplicações práticas com pesos de redes neurais,
trabalha-se com a versão discretizada por meio de histogramas com $N$
bins, recuperando a forma da \autoref{eq:shannon-discrete} a um custo
controlado de viés e variância \cite{coverthomas2006}.

Uma medida de dissimilaridade entre duas distribuições $P$ e $Q$
sobre o mesmo alfabeto é a divergência de Kullback--Leibler
\begin{equation}
  D_{\mathrm{KL}}(P\|Q) = \sum_{i=1}^{N} p_i \log_2 \frac{p_i}{q_i},
  \label{eq:kl}
\end{equation}
que é não negativa e nula se, e somente se, $P = Q$. A
\autoref{eq:kl} é central nesta dissertação porque a noção de
``desequilíbrio'' subjacente à complexidade LMC pode ser interpretada
como uma medida de afastamento entre $P$ e a distribuição de
referência (uniforme), tema retomado na \autoref{sec:lmc}.

\section{Entropia Aproximada e Entropia de Amostra}
\label{sec:apen-sampen}

O salto da teoria da informação para a análise de séries temporais
finitas exige um estimador de regularidade que opere com poucos pontos
e seja robusto a ruído moderado. A Entropia Aproximada (ApEn),
introduzida por \citeonline{pincus1991}, foi a primeira proposta a
preencher esse papel para sinais fisiológicos. Dada uma série
$\mathbf{u} = (u_1, u_2, \ldots, u_N)$, fixa-se um comprimento de
embedding $m$ e uma tolerância $r$ e constroem-se os vetores
$\mathbf{x}_i = (u_i, u_{i+1}, \ldots, u_{i+m-1})$ com
$1 \le i \le N - m + 1$. Para cada $i$ define-se
\begin{equation}
  C_i^m(r) = \frac{\#\{j : d[\mathbf{x}_i, \mathbf{x}_j] \le r\}}
                  {N - m + 1},
  \label{eq:apen-ci}
\end{equation}
onde $d[\cdot, \cdot]$ é a distância de Chebyshev (máximo das
diferenças absolutas componente a componente) e a contagem inclui o
caso $j = i$ (\emph{self-match}). A ApEn é então
\begin{equation}
  \mathrm{ApEn}(m, r, N) = \Phi^m(r) - \Phi^{m+1}(r),
  \quad
  \Phi^m(r) = \frac{1}{N - m + 1}\sum_{i=1}^{N-m+1} \ln C_i^m(r).
  \label{eq:apen}
\end{equation}

A inclusão de \emph{self-matches} em $C_i^m(r)$ é simultaneamente o
recurso que viabiliza a definição (evita logaritmo de zero) e a fonte
de seus dois principais defeitos: (i) ApEn é sistematicamente
enviesada para valores menores em séries curtas, e (ii) sua
consistência relativa --- a propriedade de que, ao se comparar duas
séries em um mesmo $r$, a ordem dos valores se preserva quando $r$
varia --- não é garantida \cite{richman2000,delgadobonal2019}. A
literatura subsequente documenta cenários em que classificações
fisiológicas baseadas em ApEn se invertem ao se mudar a tolerância,
comprometendo seu uso como marcador objetivo
\cite{yentes2013,delgadobonal2019}.

A Entropia de Amostra (SampEn), proposta por \citeonline{richman2000},
remove essas limitações por construção. Em vez de calcular o
logaritmo médio de probabilidades locais, conta-se globalmente o
número de pares de templates próximos. Sejam
\begin{align}
  B^m(r) &= \#\{(i, j) : i < j,\ d[\mathbf{x}_i, \mathbf{x}_j] \le r\},
  \label{eq:sampen-B} \\
  A^m(r) &= \#\{(i, j) : i < j,\ d[\mathbf{x}_i, \mathbf{x}_j] \le r,
                       \ |u_{i+m} - u_{j+m}| \le r\},
  \label{eq:sampen-A}
\end{align}
isto é, $B^m(r)$ conta pares de janelas de comprimento $m$ próximas e
$A^m(r)$ conta os pares dessas janelas que permanecem próximas após
um passo adicional. A SampEn é dada por
\begin{equation}
  \mathrm{SampEn}(m, r, N) = -\ln \frac{A^m(r)}{B^m(r)}.
  \label{eq:sampen}
\end{equation}
Pela construção, ambas as contagens excluem o par $(i, i)$, eliminando
o viés associado às auto-coincidências. A interpretação de
\autoref{eq:sampen} é direta e instrutiva: SampEn é o negativo do
logaritmo da probabilidade condicional de que duas janelas próximas
em comprimento $m$ permaneçam próximas em comprimento $m + 1$.
Valores baixos correspondem a séries muito regulares (a continuação
é altamente previsível); valores altos, a séries irregulares
(continuação imprevisível).

A \autoref{eq:sampen} herda independência tanto de $N$ quanto da
ordem dos pares por ser uma estatística baseada em razão de contagens
globais. Demonstrou-se ainda que, em uma ampla família de processos
estocásticos, SampEn possui consistência relativa entre tolerâncias,
propriedade que ApEn não satisfaz \cite{richman2000}. Por outro lado,
a literatura recente registra cenários particulares de processos do
tipo $\mathrm{MIX}$ contaminados por ruído em que mesmo SampEn pode
falhar essa consistência, recomendando cautela ao se interpretar
diferenças finas entre séries de naturezas distintas
\cite{delgadobonal2019}.

A escolha dos parâmetros $m$ e $r$ é prática consagrada na
literatura. Pela análise sistemática de \citeonline{yentes2013} sobre
séries curtas, recomenda-se $m = 2$ e $r$ entre $0{,}1$ e $0{,}25$ do
desvio padrão da série, com $N \ge 200$ para evitar instabilidade
estatística. Esta dissertação adota $m = 2$ e $r = 0{,}2 \times
\sigma$, com $\sigma$ tomado como desvio padrão populacional
(\texttt{ddof}=0), em consonância com a convenção do grupo de pesquisa
e com o intervalo recomendado pela literatura. Ressalva-se que tensores
pequenos --- em particular os parâmetros $\gamma$ e $\beta$ do
\emph{batch normalization} das camadas iniciais --- podem ter
comprimento inferior a $N = 200$, regime em que a estimativa de SampEn
é menos estável; esse limite é levado em conta na interpretação desses
grupos.

\section{Entropia de Amostra Bidimensional (SampEn2D)}
\label{sec:sampen2d}

A extensão bidimensional da SampEn surge da observação de que pesos
de camadas convolucionais e parâmetros de matrizes têm estrutura
espacial: a vizinhança de um peso é tão informativa quanto sua
sequência linear, e o domínio natural para análise é a matriz, não
o vetor. \citeonline{silva2014embc} apresentaram uma primeira
formulação, validada posteriormente em \citeonline{silva2016} sobre
imagens de textura sintéticas, do banco Brodatz e de cortes
histológicos de nervo sural, demonstrando que SampEn2D recupera as
propriedades qualitativas de SampEn 1D ao mesmo tempo em que
incorpora a dependência bidimensional.

Seja $\mathbf{W} \in \mathbb{R}^{H \times W}$ uma matriz de pesos
extraída de uma camada da rede. Fixe-se um tamanho de janela $m$ e
uma tolerância $r$. As janelas $m \times m$ centradas em
$(i, j)$ formam a coleção
$\mathcal{V}_m = \{ \mathbf{V}_{i,j}^{(m)} : 1 \le i \le H - m,\,
1 \le j \le W - m\}$. Para uma janela maior $(m+1) \times (m+1)$,
analogamente, escreve-se $\mathcal{V}_{m+1}$. Por construção,
restringe-se às posições $(i, j)$ tais que ambas as janelas $m$ e
$m+1$ caibam integralmente em $\mathbf{W}$, garantindo emparelhamento
entre $\mathcal{V}_m$ e $\mathcal{V}_{m+1}$. Em analogia com as
\autoref{eq:sampen-B} e \autoref{eq:sampen-A}, definem-se
\begin{align}
  B^m_{\mathrm{2D}}(r)
    &= \#\{(p, q) : p \neq q,\ d[\mathbf{V}_p^{(m)}, \mathbf{V}_q^{(m)}] \le r\},
    \label{eq:sampen2d-B} \\
  A^m_{\mathrm{2D}}(r)
    &= \#\{(p, q) : p \neq q,\ d[\mathbf{V}_p^{(m+1)}, \mathbf{V}_q^{(m+1)}] \le r\},
    \label{eq:sampen2d-A}
\end{align}
em que $p, q$ percorrem as posições válidas e $d[\cdot, \cdot]$ é a
distância de Chebyshev sobre a janela vista como vetor. A
\emph{Sample Entropy 2D} é
\begin{equation}
  \mathrm{SampEn2D}(m, r, \mathbf{W})
    = -\ln \frac{A^m_{\mathrm{2D}}(r)}{B^m_{\mathrm{2D}}(r)}.
  \label{eq:sampen2d}
\end{equation}
A interpretação probabilística replica a do caso unidimensional:
SampEn2D é o negativo do logaritmo da probabilidade de que duas
janelas $m \times m$ próximas tenham vizinhanças que permaneçam
próximas quando a janela é estendida em uma linha e uma coluna.

A literatura aponta dois aspectos práticos relevantes. Primeiro, a
SampEn2D pode produzir valores indefinidos para matrizes pequenas
quando $A^m_{\mathrm{2D}}(r) = 0$, problema que motivou propostas
alternativas baseadas em entropias de dispersão e
fuzzificação. Segundo, o custo computacional cresce
quadraticamente com o número de janelas válidas, tornando a
estimativa custosa para matrizes grandes \cite{silva2016}. Para os
fins desta dissertação, o tamanho típico das matrizes de uma camada
convolucional na ResNet-18 mantém-se confortavelmente acima do
regime patológico identificado pela literatura, e a SampEn2D é
aplicada diretamente, sem recurso a substitutos.

\section{Entropia Multiescala (MSE)}
\label{sec:mse}

A SampEn captura irregularidade em uma escala fixa. No entanto,
sistemas com correlações de longo alcance --- batimentos cardíacos
saudáveis, ruído $1/f$, dinâmica de redes complexas ---  apresentam
estrutura em múltiplas escalas temporais simultaneamente. A Entropia
Multiescala (MSE), proposta por \citeonline{costa2002} e
sistematizada em \citeonline{costa2005}, lê esse fenômeno por meio de
uma sequência de granulações grosseiras (\emph{coarse-grainings}) da
série original.

Dada uma série $\mathbf{u}$ e um fator de escala $\tau \in
\mathbb{N}^{\ast}$, constrói-se a série granulosa
\begin{equation}
  y_j^{(\tau)} = \frac{1}{\tau} \sum_{i = (j-1)\tau + 1}^{j\tau} u_i,
  \quad 1 \le j \le \lfloor N/\tau \rfloor,
  \label{eq:mse-coarse}
\end{equation}
isto é, médias de blocos não sobrepostos de comprimento $\tau$. A
escala $\tau = 1$ recupera a série original. Para cada $\tau$,
calcula-se SampEn da série granulosa, formando o perfil multiescala
$\{\mathrm{SampEn}(m, r, y^{(\tau)}) : \tau = 1, 2, \ldots,
\tau_{\max}\}$.

Um detalhe metodológico crítico: a tolerância $r$ é fixada uma única
vez a partir do desvio padrão da série \emph{original} e mantida
constante em todas as escalas, conforme prescrição de
\citeonline{costa2005}. Recomputar $r$ em cada escala invalidaria a
comparação entre escalas, pois mudanças no perfil de complexidade
ficariam confundidas com a renormalização local. A interpretação
qualitativa da curva $\mathrm{SampEn}$ versus $\tau$ separa três
classes principais: ruído branco (entropia decresce monotonicamente),
ruído $1/f$ (entropia constante ao longo das escalas) e periódicos
(entropia próxima de zero em todas as escalas). Sistemas
fisiológicos saudáveis apresentam padrão similar ao $1/f$,
enquanto patologias frequentemente colapsam para perfis
monotonicamente decrescentes \cite{costa2005}.

Esta dissertação aplica MSE 1D ao vetor de pesos achatado e, na
extensão proposta por \citeonline{silva2016} para SampEn2D, MSE 2D às
matrizes via granulação por blocos $\tau \times \tau$. Em ambos os
casos, $r$ é mantido fixo em $0{,}2 \times \sigma$ da matriz original,
com $\sigma$ calculado em população.

\section{Complexidade Estatística LMC}
\label{sec:lmc}

A entropia, isoladamente, não é uma boa medida de complexidade.
Distribuições uniformes (entropia máxima) são tão simples quanto
distribuições degeneradas (entropia mínima), no sentido de que ambas
são facilmente descritíveis e desprovidas de estrutura interna.
\citeonline{lmc1995}\,---\,doravante LMC\,---\,propuseram, como
remédio, uma medida que se anula em ambos os extremos e atinge
valores máximos em regimes intermediários de auto-organização. Para
uma distribuição $P = (p_1, \ldots, p_N)$, define-se o
\emph{desequilíbrio}
\begin{equation}
  D(P) = \sum_{i=1}^{N} \left( p_i - \frac{1}{N} \right)^{2},
  \label{eq:lmc-disequilibrium}
\end{equation}
isto é, o quadrado da distância $L_2$ até a distribuição uniforme.
A complexidade estatística LMC é o produto
\begin{equation}
  C(P) = H(P) \times D(P),
  \label{eq:lmc}
\end{equation}
em que $H$ é a entropia de Shannon, dada pela
\autoref{eq:shannon-discrete}. Como $H$ é máxima e $D$ é nula sob a
distribuição uniforme, e $H$ é nula e $D$ é máxima sob a distribuição
degenerada, $C$ se anula em ambos os extremos por construção. Para
distribuições intermediárias, $C$ assume valores estritamente
positivos, refletindo simultaneamente incerteza e afastamento do
equilíbrio \cite{lmc1995}.

A definição original gerou debate intenso. \citeonline{anteneodo1996}
demonstraram, no contexto de ensembles físicos, que $C$ na forma
\autoref{eq:lmc} \emph{não} é uma quantidade intensiva: ao se
duplicar o tamanho do sistema, $C$ não satisfaz a relação de
escala esperada de uma medida termodinâmica genuína. Esse fato
limita a transferibilidade direta de $C$ entre sistemas de
tamanhos distintos e motivou a busca por refinamentos.
\citeonline{feldman1998}, em uma revisão crítica das medidas de
complexidade estatística disponíveis na época, articularam três
desideratos: a medida deve (i) anular-se nos extremos de ordem e
desordem perfeitas, (ii) ser intensiva, e (iii) capturar a presença
de estrutura correlacional em escalas relevantes do sistema.

A resposta mais influente a essas críticas é a complexidade
estatística generalizada (SCM, \emph{Statistical Complexity
Measure}), introduzida por \citeonline{lamberti2004} e desenvolvida
por \citeonline{rosso2010}. Em vez do produto $H \times D$, define-se
\begin{equation}
  C_{\mathrm{SCM}}(P) = Q_J(P, P_e) \cdot H(P) / H_{\max},
  \label{eq:scm}
\end{equation}
em que $H_{\max} = \log_2 N$ normaliza a entropia, $P_e$ é a
distribuição de equilíbrio (tipicamente uniforme), e $Q_J(P, P_e)$ é
o desequilíbrio em divergência de Jensen-Shannon, normalizado de modo
que $0 \le Q_J \le 1$. A construção \autoref{eq:scm} é intensiva,
limitada ao intervalo $[0, 1]$ e aplicável a uma variedade de
sistemas dinâmicos contínuos e discretos.

Apesar das vantagens formais da SCM, esta dissertação adota a forma
original \autoref{eq:lmc}\,---\,$C = H \times D$ em bits, sem
normalização\,---\,por três razões metodológicas que merecem
explicitação:

\begin{enumerate}
  \item \textbf{Comparação intra-modelo, não inter-sistema.} O objeto
        de análise são as trajetórias de complexidade dos pesos de
        \emph{uma mesma rede} em \emph{um mesmo grupo de camadas} ao
        longo do treinamento. As comparações são pareadas por
        \emph{checkpoint}, \emph{seed} e dataset; o argumento de
        intensividade de \citeonline{anteneodo1996}, que diz respeito
        à escalonamento entre sistemas de tamanhos distintos, não é
        ativado nesse desenho experimental.
  \item \textbf{Sensibilidade a transições intermediárias.} A forma
        $H \times D$ é sensível justamente ao regime que interessa à
        hipótese central\,---\,a transição entre o estado inicial
        aleatório e o estado treinado. Versões normalizadas saturam
        nas extremidades do intervalo $[0, 1]$, comprimindo as
        diferenças que se busca medir. A versão não normalizada
        preserva a granularidade necessária para detectar inflexões
        sutis na trajetória.
  \item \textbf{Convenção do grupo de pesquisa.} Trabalhos anteriores
        do mesmo grupo \cite{silva2016,silva2012qsampen} adotam a
        convenção $H \times D$ em bits, sem normalização, com
        binarização adaptativa da forma $N = \min(100, \max(10,
        n/5))$, onde $n$ é o número de amostras. Manter essa
        convenção preserva a comparabilidade direta entre os
        resultados desta dissertação e os da literatura imediata.
\end{enumerate}

A normalização adaptativa do número de bins $N$ é especialmente
relevante para a aplicação a pesos de redes neurais, onde o número
de parâmetros varia ordens de grandeza entre camadas: tipicamente
poucas centenas em camadas iniciais e dezenas de milhares em
camadas profundas. A regra $N = \min(100, \max(10, n/5))$ garante
resolução suficiente para camadas pequenas sem incorrer em viés de
subamostragem em camadas grandes.

A combinação entre LMC (sensível à forma da distribuição marginal
dos pesos) e SampEn (sensível à estrutura sequencial e espacial)
fornece, portanto, dois ângulos complementares sobre o estado
interno da rede ao longo do treinamento. Os parâmetros operacionais
adotados, junto da fundamentação aqui apresentada, são retomados em
detalhe no \autoref{ch:metodologia}.

\end{claudedraft}

% ===========================================================================
\section{Estado da arte: leitura do estado interno de redes profundas}
\label{sec:trabalhos-relacionados}

Esta seção discute os trabalhos que se aproximam do problema desta dissertação.
O critério de seleção adotado foi a existência de alguma tentativa de
caracterizar o estado de generalização de uma rede a partir de seu
\emph{interior} --- pesos, representações ou dinâmica de treinamento --- em
lugar de aferi-lo exclusivamente pelo desempenho em dados rotulados. A
Seção~\ref{sec:ea-memorizacao} revisa o que se sabe sobre memorização. A
Seção~\ref{sec:ea-normas} discute a família de limites baseados em normas e
margens. A Seção~\ref{sec:ea-espectral} trata das abordagens espectrais. A
Seção~\ref{sec:ea-info} apresenta a família informação-teórica e a controvérsia
que a acompanha. A Seção~\ref{sec:ea-sintese} sintetiza e posiciona a proposta.

\subsection{Memorização e os dois regimes do aprendizado}
\label{sec:ea-memorizacao}

\citeonline{zhang2017} demonstraram que redes convolucionais de grande porte
ajustam perfeitamente rótulos atribuídos ao acaso, e que esse comportamento
persiste sob regularização explícita. O resultado é duplamente relevante para
esta dissertação: estabelece que a capacidade de memorização não é acidente, mas
propriedade das arquiteturas usuais, e fornece o protocolo de indução por
corrupção controlada de rótulos adotado no Capítulo~\ref{ch:metodologia}. A
limitação, do ponto de vista aqui assumido, é que o diagnóstico permanece
ancorado na comparação entre desempenho de treino e de teste.

\citeonline{arpit2017} refinam o quadro ao mostrar que memorização e aprendizado
de padrões não são simultâneos: redes tendem a aprender primeiro as
regularidades simples e só depois memorizar os exemplos ruidosos. A existência
de dois regimes sucessivos, e não de um processo homogêneo, é o pressuposto que
torna plausível a hipótese desta dissertação --- se há transição, pode haver
assinatura da transição. Os autores, contudo, caracterizam os regimes por
métricas que dependem dos dados, não da estrutura dos pesos.

\citeonline{power2022} descrevem o fenômeno de \emph{grokking}, em que a
generalização emerge muito depois do ajuste completo do treino, às vezes após
longo período de estagnação da validação. O fenômeno constitui contraexemplo aos
critérios ingênuos de parada antecipada: uma interrupção baseada apenas na
estagnação eliminaria justamente o regime em que a generalização acabaria por
surgir. Um indicador que leia o estado interno, e não apenas o desempenho
externo, é potencialmente capaz de distinguir estagnação de reorganização ---
questão registrada aqui como hipótese derivada, não testada neste trabalho.

\subsection{Limites de generalização por normas e margens}
\label{sec:ea-normas}

\citeonline{bartlett2017} propõem limites de generalização baseados em normas
espectrais e margens, estabelecendo que propriedades geométricas dos pesos se
relacionam com a capacidade de generalizar. \citeonline{dziugaite2017} obtêm,
por via PAC-Bayesiana, limites não vacuosos para redes sobreparametrizadas ---
resultado notável, dado que limites clássicos tornam-se vazios nesse regime.
Ambos representam avanços substantivos, mas produzem limites válidos para
\emph{um} estado da rede, tipicamente o final do treinamento, e não uma leitura
acompanhável ao longo das épocas.

\citeonline{jiang2020} conduzem o estudo empírico mais sistemático dessa
família, avaliando um grande conjunto de medidas candidatas quanto à sua relação
causal com a lacuna de generalização. A conclusão é cautelar e impõe exigência
metodológica a qualquer trabalho subsequente: muitas medidas teoricamente
motivadas correlacionam-se mal com a generalização observada, e medidas baseadas
em otimização desempenham-se melhor do que as baseadas em norma. A consequência
adotada nesta dissertação é que não basta exibir correlação entre indicador e
sobreajuste em uma execução; é preciso demonstrar \emph{discriminação} entre
execuções com e sem memorização, com controles pareados e teste de
significância.

\subsection{Abordagens espectrais}
\label{sec:ea-espectral}

\citeonline{martin2021} mostram que propriedades de cauda pesada do espectro das
matrizes de pesos permitem prever tendências de qualidade de modelos sem acesso
aos dados de treino ou de teste. Este é o trabalho mais próximo do objetivo aqui
perseguido, e convém explicitar as diferenças. A abordagem é \emph{espectral},
apoiando-se nos autovalores das matrizes, ao passo que a desta dissertação é
\emph{informação-teórica}, apoiando-se na distribuição de valores e na
organização espacial dos pesos. Além disso, o alvo em \citeonline{martin2021} é
comparar modelos já treinados, enquanto aqui o alvo é acompanhar a trajetória
temporal de um único treinamento.

\subsection{Abordagens informação-teóricas e a controvérsia do gargalo de informação}
\label{sec:ea-info}

A família mais próxima em espírito a esta dissertação é a que aplica teoria da
informação ao interior de redes profundas. \citeonline{tishby2015} propõem
analisar redes no \emph{plano de informação}, quantificando cada camada pela
informação mútua que preserva a respeito da entrada e da saída, e interpretam o
treinamento como otimização de um compromisso entre compressão e predição.
\citeonline{shwartzziv2017} desenvolvem empiricamente essa leitura e relatam
que a maior parte das épocas de treinamento é consumida não no ajuste dos
rótulos, mas na \emph{compressão} da representação, identificando duas fases
distintas na otimização --- uma de deriva rápida e outra de relaxação
estocástica.

A pertinência para o presente trabalho é evidente: trata-se de literatura que
afirma existirem transições de regime detectáveis por quantidades
informação-teóricas calculadas sobre o interior da rede, o que é precisamente a
estrutura da hipótese aqui investigada.

É imprescindível, porém, registrar que essas conclusões foram contestadas.
\citeonline{saxe2019} examinam os três principais postulados da teoria e
concluem que nenhum deles se sustenta no caso geral: a existência de duas fases
distintas, a relação causal entre compressão e boa generalização, e a atribuição
da compressão ao comportamento difusivo do gradiente estocástico. Demonstram, em
particular, que a trajetória observada no plano de informação depende
predominantemente da não linearidade empregada --- funções de saturação bilateral
produzem fase de compressão, enquanto ativações lineares ou de saturação
unilateral não a produzem --- e que parte do efeito reflete decisões de estimação
da informação mútua em redes determinísticas.

Essa controvérsia recomenda cautela e delimita o escopo do presente trabalho de
modo produtivo. As medidas aqui empregadas \emph{não} são estimativas de
informação mútua entre camadas e dados, e portanto não herdam as dificuldades de
estimação que Saxe e colaboradores identificam: LMC, SampEn e SampEn2D são
estatísticas calculadas sobre os \emph{valores dos pesos}, sem referência a
qualquer distribuição conjunta com os dados. A escolha evita o problema
metodológico central da controvérsia, ao custo de abdicar da interpretação
formal em termos de compressão de informação.

\subsection{Síntese e posicionamento da proposta}
\label{sec:ea-sintese}

\begin{table}[H]
\centering
\caption{Comparação dos trabalhos discutidos quanto à dependência de dados
rotulados, à natureza temporal da leitura e ao objeto medido.}
\label{tab:estado-arte}
\footnotesize
\begin{tabular}{llccc}
\toprule
\textbf{Trabalho} & \textbf{Objeto da medida} & \textbf{Exige rótulo} &
\textbf{Leitura temporal} & \textbf{Avisa antes} \\
\midrule
\citeonline{prechelt1998}   & Erro de validação          & Sim & Sim & Não \\
\citeonline{zhang2017}      & Desempenho                 & Sim & Não & Não \\
\citeonline{arpit2017}      & Desempenho                 & Sim & Sim & Não \\
\citeonline{bartlett2017}   & Norma espectral            & Não & Não & Não \\
\citeonline{dziugaite2017}  & Limite PAC-Bayes           & Sim & Não & Não \\
\citeonline{jiang2020}      & Diversas                   & Sim & Não & Não \\
\citeonline{martin2021}     & Espectro dos pesos         & Não & Não & Não \\
\citeonline{tishby2015}     & Informação mútua           & Sim & Não & Não \\
\citeonline{shwartzziv2017} & Informação mútua           & Sim & Sim & Não \\
\textbf{Esta dissertação}   & Complexidade dos pesos     & Não & Sim & Sim \\
\bottomrule
\end{tabular}
\fonte{Elaborado pelo autor.}
\end{table}

A Tabela~\ref{tab:estado-arte} organiza os trabalhos segundo três eixos. Quanto
à dependência de rótulos, apenas \citeonline{bartlett2017} e
\citeonline{martin2021} dispensam-nos, e ambos por via espectral. Quanto à
natureza temporal, \citeonline{arpit2017} e \citeonline{shwartzziv2017}
acompanham a trajetória do treinamento, mas ambos requerem acesso aos dados ---
o primeiro ao desempenho, o segundo à distribuição conjunta necessária para
estimar informação mútua. A intersecção entre \emph{dispensar rótulos} e
\emph{acompanhar a trajetória} permanece, até onde a revisão alcançou, não
ocupada.

A proposta desta dissertação situa-se nessa intersecção e diferencia-se da
literatura em quatro aspectos: (i) a medida é informação-teórica, mas incide
sobre a distribuição de valores e a organização espacial dos pesos, não sobre o
espectro das matrizes nem sobre informação mútua com os dados; (ii) a leitura é
temporal, produzindo uma trajetória por época em lugar de um escalar ao final do
treino; (iii) nenhum rótulo é consumido, o que a torna aplicável em cenários de
anotação escassa; e (iv) a validação é feita por controles estritamente pareados
com indução explícita de memorização, em resposta direta à advertência de
\citeonline{jiang2020} quanto a correlações espúrias.

Registre-se, por fim, uma limitação da própria revisão. A literatura de medidas
de generalização é extensa e cresce rapidamente; a seleção aqui apresentada
privilegiou trabalhos fundacionais de alta citação e os mais diretamente
relacionados ao problema, não constituindo revisão sistemática no sentido
metodológico do termo.
